\documentclass[12pt]{article}

\usepackage[a4paper,margin=24mm]{geometry}
\usepackage{amsmath,amssymb,bm}
\usepackage{booktabs}
\usepackage{graphicx}
\usepackage{microtype}
\usepackage[numbers,sort&compress]{natbib}
\usepackage[hidelinks]{hyperref}
\usepackage{xcolor}

% The generated result macros retain historical digit-bearing names. TeX normally
% terminates a control word before a digit, so temporarily classify 1 and 5 as
% letters while reading the generated file and access the names through
% \csname below. This keeps the frozen reporting output reproducible.
\def\restoreVfifteendigitcatcodes{\catcode`1=12\catcode`5=12}
\catcode`5=11
\catcode`1=11
% Generated from V15 aggregate results; do not edit manually.
\newcommand{\V15Trajectories}{48}
\newcommand{\V15Decisions}{34,560}
\newcommand{\V15ClustersPerModel}{6}
\newcommand{\V15GPUHours}{49.41}
\newcommand{\V15OriginalGPUHours}{19.19}
\newcommand{\V15ReconstructionGPUHours}{49.41}
\newcommand{\V15OrphanFormalDecisions}{197}
\newcommand{\V15OrphanFormalGPUHours}{0.328}
\newcommand{\V15UnrecordedInfrastructureCalls}{1}
\newcommand{\V15HOne}{42.263}
\newcommand{\V15HOneCI}{24.316 to 59.303}
\newcommand{\V15HTwo}{0.05438}
\newcommand{\V15HTwoCI}{0.03459 to 0.07548}
\newcommand{\V15HThree}{0.04845}
\newcommand{\V15HThreeCI}{0.02190 to 0.07526}
\newcommand{\V15HFour}{52.541}
\newcommand{\V15HFourCI}{35.406 to 68.673}
\newcommand{\V15HOneDisposition}{supported}
\newcommand{\V15HTwoDisposition}{supported}
\newcommand{\V15HThreeDisposition}{supported}
\newcommand{\V15HFourDisposition}{supported}
\newcommand{\V15QwenCorrQuench}{0.00238}
\newcommand{\V15QwenCorrQuenchCI}{-0.01750 to 0.02277}
\newcommand{\V15GraniteCorrQuench}{-0.02032}
\newcommand{\V15GraniteCorrQuenchCI}{-0.05486 to 0.00947}
\newcommand{\V15QwenTauMemory}{-2.46995}
\newcommand{\V15QwenTauMemoryCI}{-5.79214 to 0.17724}
\newcommand{\V15GraniteTauMemory}{-4.14109}
\newcommand{\V15GraniteTauMemoryCI}{-19.88286 to 11.60067}
\newcommand{\V15QwenBinderQuench}{-4.80504}
\newcommand{\V15QwenBinderQuenchCI}{-7.42824 to -2.20832}
\newcommand{\V15GraniteBinderQuench}{0.16643}
\newcommand{\V15GraniteBinderQuenchCI}{-0.09447 to 0.52063}

\restoreVfifteendigitcatcodes
\newcommand{\E}{\mathbb E}
\newcommand{\KL}{D_{\mathrm{KL}}}
\newcommand{\figroot}{figures}
\newcommand{\resultfigure}[4]{%
\begin{figure}[htbp]\centering
\IfFileExists{\figroot/#1}{\includegraphics[width=#2]{\figroot/#1}}{\fbox{\parbox[c][45mm][c]{0.88\textwidth}{\centering Figure generated after formal analysis:\\[2mm]\texttt{\detokenize{#1}}}}}
\caption{#3}\label{#4}\end{figure}}

\title{Statistical-mechanical characterization of memory and quench response in state-separated LLM-agent networks}
\author{Anonymous working manuscript}
\date{23 August 2026}

\begin{document}
\maketitle

\begin{abstract}
We treat a network of state-separated, locally informed large-language-model
(LLM) agent instances as a finite interacting stochastic system.  Each
persistent identity owns a categorical belief, a committed action, bounded
private memory, a local field, an inbox and outbox, and a typed local action
interface;
a random-sequential scheduler offers update opportunities but does not choose
the resulting state.  We distinguish the complete augmented simulator state,
its observable belief--action projection, and a rolling macroscopic
representation.  This separation exposes how omitted private memory can act
as a hidden historical coordinate.  We report collective order, uncertainty,
dependence, an effective symmetric-layer energy, fluctuations, and
coarse-grained path-reversal divergence.  We first correct a derived-analysis
error in an earlier Qwen quench experiment without changing any trajectory.
We then conduct a prospectively specified two-family experiment with
\csname V15Trajectories\endcsname{} trajectories,
\csname V15Decisions\endcsname{} local decisions, and six
independent graph/environment clusters per model.  A field reversal and
restoration is paired with nominal evolution, genuine persistent memory, and a
format-matched scrambled-history control.  The independent-family quench
effect is \csname V15HOne\endcsname{} (95\% interval
\csname V15HOneCI\endcsname) within-model distance units.  Genuine
memory changes bias-adjusted reversal divergence relative to Markovized and
scrambled-history controls by \csname V15HTwo\endcsname{}
(\csname V15HTwoCI\endcsname) and \csname V15HThree\endcsname{}
(\csname V15HThreeCI\endcsname) nats per attempted update.  The fixed
early-to-late restoration contrast is \csname V15HFour\endcsname{}
(\csname V15HFourCI\endcsname) distance units.  All four effects meet their
prespecified criteria, although the memory contrasts are larger and more uniform in
Granite than in Qwen.  An out-of-sample kinetic
surrogate identifies which response components admit a simple local closure
and which remain model- and history-dependent.  These finite-size results
support a reduced statistical-mechanical description without interpreting
decoding temperature as physical temperature, reference energy as literal
energy, or projected path divergence as exact thermodynamic entropy
production.
\end{abstract}

\noindent\textbf{Keywords:} statistical mechanics of agents; language-model
agents; quench dynamics; entropy rate; multi-information; temporal asymmetry.

\section{Introduction}
\label{sec:introduction}

Large-language models are increasingly studied not only as isolated
prompt--response systems but as components of interacting agent populations
\cite{wang2024agents,guo2024multiagents}.  Once agents exchange messages,
retain histories, and act through a network, an output from one local decision
can become part of another agent's next input.  This feedback permits
amplification, cascades, coordination, fragmentation, lock-in, delayed
response, and incomplete recovery even when every component uses fixed model
weights.  The capability of one model instance therefore does not determine
the trajectory of the group, because repeated circulation can create modes of
organization that no single generation exhibits.  The scientific question
becomes what dynamical organization emerges when many stochastic
language-model instances interact
over time, a question that matters for the analysis, design, reliability, and
control of multi-agent AI systems.

At the component level, an LLM agent places a generative model inside a
recurrent loop of observation, decision, action, and environmental feedback.
Its future inputs can contain its own earlier outputs, received messages, and
changes caused by earlier actions.  Population trajectories can consequently
depend on message content, network position, update order, retained context,
and decoding randomness.  Conventional benchmarks can characterize useful
local capabilities, but they do not reveal whether network feedback produces
stable organization, amplifies a perturbation, stores hidden history, or
returns the population to its earlier regime.  Those are properties of the
interacting process rather than of any response considered alone.

Recent agent architectures establish both the feasibility of such processes
and the source of their complexity.  Role-playing agents cooperate through
natural-language exchange \cite{li2023camel}, while generative agents combine
observation, planning, reflection, and persistent natural-language memory in a
simulated social environment \cite{park2023generative}.  These systems show
that retrieved history and communication can materially condition later
actions.  Yet architecture and application studies are commonly evaluated by
task completion or behavioral plausibility.  Those objectives do not by
themselves specify a controlled stochastic process, distinguish private state
from observed state, or identify which population coordinates retain
dynamical information.

Computational social simulation and population experiments make the same
measurement problem visible from another direction.  Conditioned language
models reproduce portions of survey-response distributions
\cite{argyle2023}, and simulated samples replicate some established behavioral
experiments while exhibiting model-specific distortions \cite{aher2023}; the
broader opportunities and validity limits of LLM-empowered agent-based
simulation are reviewed in \cite{gao2024abm}.  Networked agents have generated
consensus and fragmentation in opinion experiments \cite{chuang2024}, shared
naming conventions and collective bias \cite{flintashery2025}, and
model-dependent cooperation and coordination in repeated games
\cite{akata2025}.  Together, these studies establish that interaction can
produce reproducible group-level outcomes, while cautioning against treating
semantic plausibility as a dynamical explanation or an LLM population as a
faithful proxy for human behavior.

The interaction graph makes collective organization irreducibly networked.
Who can influence whom determines the paths along which information and
perturbations propagate, while topology, local versus collective information,
and community structure can alter the resulting dynamics.  Statistical
mechanics and network science provide tools for connecting such interaction
structure to collective behavior \cite{albert2002,boccaletti2006}; studies of
cooperative processes on networks also emphasize that topology and finite
size can qualitatively affect observed regimes \cite{dorogovtsev2008}.  A
population mean suppresses this organization: identical mean order can
coexist with different fluctuations, graph-distance correlations, community
patterns, and routes of influence.  Graph distance supplies a natural
coordinate for asking whether local perturbations remain localized or become
collectively organized.  Network-aware observables are therefore
needed even when the scientific target is a population-level response.

More broadly, statistical mechanics supplies a disciplined micro-to-macro
measurement framework.  Agent-based models ask how local rules and
heterogeneous interactions generate system-level behavior
\cite{bonabeau2002}, and statistical physics has developed quantitative
accounts of opinion, cultural, language, and consensus dynamics
\cite{castellano2009,baronchelli2018}.  Canonical stochastic models show how
local updates generate collective relaxation \cite{glauber1963}, how a
potential can constrain stationary interaction dynamics \cite{blume1993}, and
how nonreciprocity can support different collective structures
\cite{fruchart2021}.  Local language-conditioned transitions are much more
complicated than spin updates, but population order, dependence, correlation,
response, and distributions can still be measured.  For a finite population
observed over finite windows, fluctuations and variation across independent
realizations carry information that a mean trajectory discards.  This use of
statistical mechanics does not require literal particles, a thermodynamic
limit, or a phase transition.

Controlled perturbation is important because nominal or steady observations
underdetermine dynamics.  Similar baseline averages could describe a robustly
organized system, a fragile one, a system that responds strongly and
recovers, or one with delayed or incomplete recovery.  Reversing an external
field and then restoring it probes finite-horizon sensitivity, response,
relaxation, and recovery directly.  Comparing that trajectory with a matched
nominal arm also separates intervention-associated change from ordinary
evolution over the same observation window.  A quench is used here as a
dynamical intervention, not as evidence that the network is an equilibrium
material.

Persistent history creates a distinct coarse-graining problem.  The complete
agent process may include private textual state that is absent from a measured
belief--action projection, so the observed population trajectory need not be
Markov even when the augmented simulator state is.  Comparing persistent and
Markovized prompts tests the effect of retaining history; comparing genuine
history with a scrambled, format-matched control helps separate temporally
related information from the mere presence of additional prompt material.
Temporal correlation then measures persistence, while path reversal tests
ordering in the projected trajectory.  State-space currents and trajectory
reversal have exact thermodynamic interpretations in appropriate fully
observed Markov systems \cite{seifert2012,roldan2012}, but coarse-graining can
hide dissipation and make inferred bounds estimator-dependent
\cite{kawaguchi2013,busiello2019}.  The path statistic used here is therefore
projected temporal asymmetry, not literal entropy production.

Recent LLM-population studies increasingly approach these questions with
statistical-physics concepts: group size can alter collective misalignment
\cite{flint2025}, topology and self--social weighting can change conformity
\cite{han2026}, intrinsic fields can outweigh fitted neighbor coupling
\cite{denobili2026}, debate noise can produce model-dependent finite-size
crossovers \cite{okawa2026}, and high-dimensional interaction data can admit
compact dynamical regimes \cite{el2026}.  This progress narrows the gap, but
the joint problem remains insufficiently resolved: controlled state and
information boundaries, field reversal and restoration, genuine and placebo
history interventions, spatial and temporal diagnostics, cluster-level
inference, and an out-of-sample closure test are rarely examined together.
A task score or single consensus coordinate cannot answer all of these
questions.

The present design combines measurements because each resolves a different
ambiguity.  Belief and action order, overlap, effective compatibility,
fluctuations, entropy, and dependence characterize alignment and synchronous
organization.  Connected graph-distance correlation detects spatial
reorganization that mean order can miss; finite-window autocorrelation and
Binder statistics describe persistence and order-parameter shape without
establishing equilibrium correlation times, critical slowing down, or a phase
transition.  Coarse-grained block reversal compares temporal ordering across
Markovized, persistent, and scrambled-history conditions.  Matched quench
controls and exact inference over independent graph--environment clusters
define the evidential unit.  Finally, an out-of-sample kinetic surrogate asks
whether a small set of macroscopic coordinates can reproduce aspects of the
full response and, equally importantly, where that reduced closure fails.

The contribution is not a general-purpose agent framework, a performance
benchmark, or a model of human society.  It is a controlled finite-size study
that treats interacting LLM instances as measurable stochastic components and
examines their collective response to external fields, retained and control
histories, network organization, and coarse-graining.  The direct experiments
use two pinned instruction-model families, $N=16$ reciprocal modular networks,
and finite observation windows; the interpretation is effective and
model-specific, with no thermodynamic-limit, exact physical
entropy-production, or universality claim.  Section~\ref{sec:system-protocol}
defines the network and interventions, Section~\ref{sec:observables-inference}
defines measurement and inference, and Section~\ref{sec:evidence-hierarchy}
sets out the study roles.  Sections~\ref{sec:quench-results}--\ref{sec:memory-results}
present the quench and memory evidence, Section~\ref{sec:closure} tests reduced
closure, and Sections~\ref{sec:discussion}--\ref{sec:conclusions} give the
interpretation and conclusions.

\section{State-separated LLM-agent networks and quench protocols}
\label{sec:system-protocol}

\subsection{Network state and update dynamics}
\label{sec:state}

For agent $i$ at attempted update $t$, the primary categorical coordinates are
belief $b_i(t)\in\{-1,+1\}$ and committed action
$a_i(t)\in\{-1,+1\}$.  Its bounded local state also includes confidence,
commitment status, workload, a memory code, at most three private history
entries, an inbox, and an outbox.  Display labels are counterbalanced and
decoded back to latent spins.  The microscopic projection is
\begin{equation}
Y_t=(\bm b_t,\bm a_t,\bm c_t,\bm m_t),
\end{equation}
but this is not the complete simulator state.

Let
\begin{equation}
\Xi_t=(S_{1,t},\ldots,S_{N,t},G_t,{\cal D}_t,q_t,R_t)
\end{equation}
denote all agent-private states $S_{i,t}$, graph state $G_t$, delivery operator
${\cal D}_t$, quench phase $q_t$, and the explicitly specified randomness
state $R_t$.  With fixed model weights and tokenizer, reconstructed local
prompts, and a defined scheduler, $\Xi_t$ induces a time-inhomogeneous
transition kernel under the external quench protocol.  The evaluator's map
$Y_t=\phi(\Xi_t)$ deliberately omits some prompt history and randomness.
Consequently $Y_t$ and any rolling macrostate need not be Markov.

Figure~\ref{fig:architecture} summarizes the information boundaries and the augmented-to-macroscopic map.

\resultfigure{figure01_architecture.pdf}{0.96\textwidth}{%
Agent information boundaries and the augmented-to-macroscopic map.  The
complete state $\Xi_t$ includes private memory, graph and schedule state; the
evaluator observes $Y_t=\phi(\Xi_t)$ and constructs rolling
$Z_t=\psi(Y_{t-w+1:t})$.  The scheduler offers a turn and delivers a packet but
does not supply the model's belief, action, memory, or signal.}{fig:architecture}

At each attempted update, a random permutation gives every agent exactly one
opportunity per sweep.  Only the scheduled instance receives its private
observation, current local state, currently delivered inbox, and, in a memory
condition, its own bounded history.  It returns a validated typed object
containing belief, action, confidence, commitment, memory state, outgoing
signal, and tool action.  One bounded repair may correct syntax; an invalid
result causes no centrally selected scientific action.  A serialized packet
is sent on one graph edge.  Fingerprints of nonrecipient peer-private states
are compared around every transition.

The corresponding information-boundary and implementation checks are described in Appendix~\ref{app:provenance}.

\subsection{History conditions and matched controls}
\label{sec:history-conditions}

The Markovized prompt is reconstructed from the explicitly recorded current
state and does not carry free conversational history.  The persistent prompt
adds the instance's last three own entries in the fixed format
\texttt{time, belief, action, memory}.  Thus two trajectories with identical
visible $(\bm b_t,\bm a_t)$ can have distinct conditional future laws:
\begin{equation}
P(Y_{t+1}\mid Y_t,M_t)\ne P(Y_{t+1}\mid Y_t,M'_t).
\end{equation}
If $M_t$ is projected out, conditional dependence on older visible states can
remain.  This is a concrete mechanism for increased block reversal divergence,
not a proof of physical dissipation.

The scrambled-history control isolates prompt presence and approximate
length.  It uses only the same agent's past opportunity times, while belief,
action, and memory content are deterministically randomized from a prespecified
control seed.  No donor agent, future time, or hidden environmental outcome is
used.  Genuine and scrambled entries have the same section, count, field
names, and format.  Token counts and their paired differences are measured;
turn-by-turn equality is not assumed.

Prompt-balance diagnostics for the scrambled-history control are shown in Figure~\ref{fig:prompt-balance} of Appendix~\ref{app:sensitivities}.

\subsection{Field quench and restoration protocol}
\label{sec:quench-protocol}

Each trajectory has 15 baseline sweeps, 15 intervention sweeps, and 15
restoration sweeps.  In the field-quench arm the vector of balanced private
fields changes sign at the first boundary and returns to its initial value at
the second.  Nominal trajectories retain the initial field.  This follows the
logic of a finite-system quench protocol, while making no quantum or
thermodynamic-limit analogy.  The initial
change is the quench; restoration is the counter-quench.

Each family contributes six new graph/environment clusters.  Each cluster has
$N=16$, a reciprocal modular fixed-degree graph, coupling $J=0.8$, balanced
disordered initialization, and four matched trajectories: nominal Markovized,
field-quench Markovized, field-quench persistent memory, and field-quench
scrambled history.  Graphs, fields, initial states, label maps, update orders,
recipient variates, and inference seeds are paired where technically
meaningful.

\section{Statistical-mechanical observables and\texorpdfstring{\\}{ }inferential design}
\label{sec:observables-inference}

\subsection{Macroscopic and information-theoretic observables}
\label{sec:order-compatibility}

Belief and action order and their overlap are
\begin{equation}
m_b=\frac{1}{N}\sum_i b_i,\qquad
m_a=\frac{1}{N}\sum_i a_i,\qquad
q_{ba}=\frac{1}{N}\sum_i b_i a_i.
\end{equation}
Disagreement is the fraction of connected pairs with unequal states.  For a
symmetric comparison layer $W_s=(W+W^\top)/2$, we define
\begin{equation}
H_{\mathrm{ref}}=-\frac{J_b}{2}\bm b^\top W_s\bm b
-\frac{J_a}{2}\bm a^\top D_s\bm a
-K\bm a^\top\bm b-\bm h^\top\bm b-\bm g^\top\bm a .
\label{eq:href}
\end{equation}
We report $e_{\mathrm{ref}}=H_{\mathrm{ref}}/N$.  Equation~\eqref{eq:href} is
an effective compatibility coordinate anchored in a
symmetric reference model.  We do not infer a Gibbs law for the LLM process
and do not call this literal energy.  Likewise
$N\operatorname{Var}(e_{\mathrm{ref}})$ is an energy-fluctuation observable,
not heat capacity unless an independently justified effective temperature
exists.

The belief susceptibility is reported without physical-temperature
normalization,
\begin{equation}
\chi_b=N\operatorname{Var}(m_b),
\end{equation}
within a prescribed rolling window.

For discretized collective words $x$ in a rolling window, Shannon
configuration
entropy is
\begin{equation}
S_{\mathrm{config}}=-\sum_x \hat p(x)\log\hat p(x).
\end{equation}
This is the usual plug-in Shannon functional \cite{shannon1948}.
The entropy rate is a bounded-history plug-in estimate of next-word
unpredictability.  Single-variable marginal entropy measures local occupancy;
pairwise mutual information measures dependence between variables.  Total
correlation \cite{watanabe1960},
\begin{equation}
{\cal T}=\sum_{k=1}^{2N}H(X_k)-H(X_1,\ldots,X_{2N}),
\end{equation}
separates marginal uncertainty from synchronous joint dependence.

The plug-in joint entropy of 32 binary variables in a short window has a large
finite-sample floor \cite{paninski2003}.  We therefore retain the raw estimate
but independently circular-shift each variable 200 times.  Shifts preserve
each marginal occupancy and within-variable temporal ordering while breaking
synchronous cross-variable relations.  We report raw, null mean, their
untruncated difference, and a normalized difference divided by the sum of
marginal entropies.  The same audit is applied to pairwise and edge mutual
information.  A negative adjusted estimate remains negative.

The primary five-sweep representation is
\begin{align}
Z_t=(&m_b,m_a,q_{ba},e_{\mathrm{ref}},v_E,S_{\mathrm{config}},h_{\mathrm{rate}},
{\cal T}_{\mathrm{adj}},I_{\mathrm{pair,adj}},I_{\mathrm{edge,adj}},\nonumber\\
&\chi_b,\rho_b,D_b)_t .
\end{align}
Three- and seven-sweep windows are prespecified sensitivities.  This is an
observable reduced representation, not a complete thermodynamic state.

\subsection{Temporal asymmetry, spatial correlation, and persistence diagnostics}
\label{sec:path-asymmetry}

For a discrete observed word sequence and block length $L$, let
$\hat P_L(x_0,\ldots,x_L)$ be the regularized empirical forward block law.
We compute
\begin{equation}
{\cal I}_L=\frac{1}{L}
\KL\left[\hat P_L(x_0,\ldots,x_L)\parallel
\hat P_L(x_L,\ldots,x_0)\right]
\label{eq:pathkl}
\end{equation}
in nats per attempted update.  The primary $L=3$ estimator uses pseudocount
0.5.  Values at $L\in\{2,4\}$ and pseudocounts 0.1 and 1.0 are sensitivities.
Five hundred time permutations provide a panel-specific finite-length floor,
and the adjusted statistic is the observed value in
equation~\eqref{eq:pathkl} minus its mean floor.
Matched arms share the same shuffle tape.  The machine-readable sensitivity
table retains both the empirical permutation-null interval and the Monte Carlo
standard error of its mean.  These quantify different objects: dispersion of
the null statistic and numerical uncertainty in the estimated bias floor,
respectively.  The added uncertainty columns reproduce, rather than replace,
the prespecified floor mean and adjusted contrast.

For a fully observed stationary Markov chain with transition kernel $P$ and
stationary law $\pi$, the current
\begin{equation}
J_{xy}=\pi_xP_{xy}-\pi_yP_{yx}
\end{equation}
can define an exact Markov entropy-production rate under the usual support
conditions.  We do not apply that interpretation to projected LLM trajectories.
Persistent histories, semantic content, and bounded observations can violate
first-order closure.  Equation~\eqref{eq:pathkl} is instead a coarse-grained
time-reversal diagnostic and, under additional assumptions, may be a lower
bound on hidden irreversibility \cite{roldan2012,busiello2019,liang2023}.

Three secondary observables resolve
spatial organization, temporal persistence, and the shape of the finite-size
order-parameter distribution.  First, for a phase window $W$ we define the
connected belief correlation at graph distance $d$,
\begin{equation}
C_b(d;W)=\frac{1}{|{\cal P}_d|}\sum_{(i,j)\in{\cal P}_d}
\left[\langle b_i b_j\rangle_W-\langle b_i\rangle_W
\,\langle b_j\rangle_W\right],\qquad
{\cal P}_d=\{(i,j):i<j,\ d_G(i,j)=d\}.
\label{eq:connected-correlation}
\end{equation}
Shortest paths use the actual reciprocal modular delivery support.  The
bracket averages post-update observations at attempted-update resolution
within one phase of one
trajectory; node pairs are averaged inside that trajectory and never counted
as independent replicates.

Second, the normalized magnetization autocorrelation and its fixed truncation
are
\begin{align}
\rho_m(\ell)&=\frac{\langle[m_b(t)-\bar m_b]
\,[m_b(t+\ell)-\bar m_b]\rangle}
{\langle[m_b(t)-\bar m_b]^2\rangle},\nonumber\\
\tau_{\mathrm{int}}^{(\ell_{\max})}&=\frac{1}{2}+
\sum_{\ell=1}^{\ell_{\max}}\rho_m(\ell).
\label{eq:autocorrelation}
\end{align}
The primary descriptive truncation is $\ell_{\max}=2N=32$ attempted updates;
$N$ and $3N$ are fixed sensitivities.  Zero-variance phase series are left
undefined rather than assigned artificial persistence.  The quantity is a
finite-window correlation sum reported in attempted-update units, not evidence
of critical slowing down or an effective sample-size estimate.  Its absolute
short-lag scale includes persistence induced by changing at most one agent on
each random-sequential update; matched arms share that update convention.
Quench and restoration windows can also be nonstationary, so their sums are
finite-window persistence diagnostics rather than equilibrium integrated
correlation times.

Third, we calculate the Binder cumulant
\begin{equation}
U_4=1-\frac{\langle m_b^4\rangle}
{3\langle m_b^2\rangle^2}
\label{eq:binder}
\end{equation}
separately for each phase and trajectory \cite{binder1981}.  Near-zero second moments are
reported as undefined.  With only $N=16$, equation~\eqref{eq:binder} is a
finite-size distribution-shape diagnostic; no crossing, critical point, or
thermodynamic-limit transition is inferred.  A post-reconstruction estimator
sensitivity repeats the calculation over each full phase and its early and
late halves.  It compares the mean of the six trajectory-level cumulants with
a pooled-moment construction; both uncertainty calculations resample complete
graph--environment clusters, never updates.

\subsection{Recovery estimands and exact inference}
\label{sec:inferential-design}

For each held-out graph/environment cluster and model, the nominal manifold is
fitted only from nominal trajectories in the other five clusters.  Training
medians impute nonfinite coordinates; training means and scales standardize
features; shrinkage covariance with a prespecified ridge fraction defines a
Mahalanobis-like distance.  The nominal 95th-percentile threshold is computed
from those training distances.  The held-out cluster contributes neither a
scale nor a threshold.

We summarize peak departure, integrated response, path length, final residual,
and threshold re-entry.  The prospective cross-model recovery contrast is
\begin{equation}
\Delta d_{\mathrm{rec}}=
\overline d_{31:35}-\overline d_{41:45},
\end{equation}
which can be positive, zero, or negative.  It replaces a historical maximum
minus final-window quantity whose sign was structurally constrained.

A complete graph/environment trajectory cluster is the
inferential unit.

The hypotheses were specified before the cross-model outcomes were examined.
H1 compares Granite field-quench peak
departure with matched nominal evolution at one-sided $\alpha=0.02$.  H2
compares persistent with Markovized adjusted reversal divergence.  H3 compares
persistent with scrambled history.  H4 tests the fixed recovery contrast.
H2--H4 form a one-sided Holm family with total $\alpha=0.03$.  Exact cluster
sign flips are the primary tests; deterministic 10,000 cluster bootstraps give
intervals.  The sign-flip enumeration is exact conditional on the observed
paired magnitudes under a sign-symmetry null; it is not described as treatment-
label randomization.  Agent and time-point counts do not increase the
inferential sample size.

The exact-test, multiplicity, and cluster-bootstrap details are retained in Appendix~\ref{app:study-inference}.

\section{Study design and evidential roles}
\label{sec:evidence-hierarchy}

\subsection{Exploratory discovery and corrected earlier evidence}
\label{sec:memory-study-roles}

The memory evidence developed in three stages: an exploratory experiment, a
subsequent prospective replication, and the present cross-model extension with
an additional history control.  Figure~\ref{fig:memory} displays these stages
separately so that the earlier estimates are not retroactively pooled into the
cross-model confirmatory interval.

\resultfigure{figure02_memory_evidence_stages.pdf}{0.94\textwidth}{%
Exploratory discovery, prospective replication, and cross-model memory contrasts are
shown separately.  Points are graph/environment-cluster effects; intervals
resample complete clusters.  The exploratory and replication-stage estimates
are not retroactively pooled into the cross-model confirmatory interval.}{fig:memory}

The audit of the earlier quench experiment preserves every recorded prompt,
completion, categorical choice, and trajectory.  The original analysis
computed a field panel's recovery
threshold from its own baseline, despite protocol text specifying a
training-only threshold.  The corrected analysis passes an explicit
cluster-to-training-threshold map from leave-one-cluster-out fitting into the
recovery summaries.

The earlier study's third hypothesis used the maximum recovery distance minus
the final-five mean.  That value remains numerically reproducible, but the
construction is
nonnegative for nearly every nonconstant trajectory.  Its raw and Holm-adjusted
sign-flip probabilities are retained as historical audit values and removed
from inferential support.  The corrected report distinguishes
whether the numerical criterion was met, whether a directional test is valid,
and whether the time-resolved path is consistent with return.

Figure~\ref{fig:v14series} displays the corrected time series from the earlier quench experiment.

\resultfigure{figure03_corrected_quench_time_series.pdf}{0.94\textwidth}{%
Corrected time series from the earlier experiment for nominal and field-reversal trajectories.  The
yellow interval marks the quench; the dotted boundary marks restoration.
Reference energy is effective, entropy is configuration diversity, and
distance is measured in the cluster-excluded nominal geometry.}{fig:v14series}

Cluster-level recovery and delayed-audit material is retained in Appendix~\ref{app:study-inference}, including Figures~\ref{fig:v14recovery} and~\ref{fig:v14audit}.

\subsection{Cross-model confirmation and control logic}
\label{sec:cross-model-design}

The original model family is represented by Qwen/Qwen2.5-7B-Instruct at immutable revision
{\footnotesize\texttt{a09a3545\allowbreak 8c702b33\allowbreak eeacc393\allowbreak d1030632\allowbreak 34e8bc28}}.  The independent family is
IBM Granite-3.3-8B-Instruct at revision
{\footnotesize\texttt{51dd4bc2\allowbreak ade4059a\allowbreak 6bd87649\allowbreak d68aa11e\allowbreak 4fb2529b}} \cite{qwen2024,granite33model}.
Both use the same pinned inference stack, NF4 double quantization, BF16
computation, decoding temperature 0.5, top-$p=0.9$, and at most 96 generated
tokens.  The decoding temperature is a model-generation setting, not an
inferred thermodynamic temperature.

The pre-outcome engineering pilot was restricted to loading, schema validity, occupancy,
both transition directions, timing, information isolation, delivery, prompt-token balance,
latency, and projected compute.  It does not calculate a quench, memory, or
irreversibility contrast.  The prospective design contains
\csname V15Trajectories\endcsname{} trajectories and
\csname V15Decisions\endcsname{} attempted local decisions.  Complete
unfavorable trajectories are retained.

\section{Quench response, restoration, and finite-size organization}
\label{sec:quench-results}

\subsection{Cross-model response and recovery}
\label{sec:field-response}

The Granite-family H1 contrast is \csname V15HOne\endcsname{} distance units
(95\% interval \csname V15HOneCI\endcsname); its prespecified disposition is
\csname V15HOneDisposition\endcsname{}.  Figure
\ref{fig:crossquench} displays model-specific distances rather than pooling
their absolute scales.  A positive contrast indicates a field-driven
departure beyond nominal variation; a nonpositive or uncertain contrast
would delimit cross-family replication.  The complete cluster table is
retained regardless of direction.

All six Granite cluster effects are positive.  The separately retained Qwen
contrast is also positive in all six clusters, with a mean field-minus-nominal
peak of 112.691 units.  These absolute values are not compared across models:
each leave-one-cluster-out geometry is standardized and fitted within model.

\resultfigure{figure04_cross_model_quench.pdf}{0.94\textwidth}{%
Model-specific field-quench and nominal trajectories.  Curves are means over
six graph/environment clusters and bands show between-cluster dispersion.
Nominal fits, scaling, covariance, and thresholds exclude the displayed
cluster.}{fig:crossquench}

Restoration H4 is \csname V15HFour\endcsname{} distance units (95\% interval
\csname V15HFourCI\endcsname), with prespecified disposition
\csname V15HFourDisposition\endcsname{}.  Unlike the earlier experiment's
historical recovery statistic, this fixed-window quantity permits either sign.  Recovery time and
threshold crossing remain trajectory descriptors, not an equilibrium
relaxation theorem.  The quench and counter-quench paths need not retrace one
another because messages, actions, and history alter the state reached at
restoration.

The fixed-window decline occurs in every model--cluster pair, but it does not
imply complete return within the observation horizon.  Every Qwen trajectory
re-enters its training-nominal threshold six sweeps after restoration.  Only
one of six Granite trajectories re-enters by sweep 45; the other five retain
final-five means above their own thresholds.  Thus H4 supports motion toward
the restored nominal regime, while Granite exhibits slower or incomplete
finite-horizon recovery.

\subsection{Spatial and finite-window organization}
\label{sec:spatial-results}

Equations~\eqref{eq:connected-correlation}--\eqref{eq:binder} are secondary
descriptive analyses of the same recorded trajectories; they neither modify the
confirmatory estimands nor introduce additional simulation arms.  They were
specified during reconstruction after the confirmatory aggregate results were
known and are therefore not a new prospective hypothesis family.  Connected
correlation matrices are first calculated within every complete
model--graph--environment trajectory and phase.  Their graph-distance profiles
are then summarized over the six independent clusters.  Figure
\ref{fig:correlations} shows cluster-averaged matrices for the persistent-
history quench window and phase-resolved distance profiles.  Community
boundaries are fixed by the graph construction rather than selected from the
observed correlations.

At graph distance one, the persistent-history quench-minus-baseline connected
correlation changes are \csname V15QwenCorrQuench\endcsname{}
(\csname V15QwenCorrQuenchCI\endcsname) for Qwen and
\csname V15GraniteCorrQuench\endcsname{}
(\csname V15GraniteCorrQuenchCI\endcsname) for Granite.  These cluster-level
descriptive intervals quantify finite-window reorganization; their signs are
not assigned a confirmatory disposition.

\resultfigure{figure05_graph_distance_correlations.pdf}{0.94\textwidth}{%
Connected belief correlations in the persistent-history arm.  The upper
panels average complete $16\times16$ connected-correlation matrices over all
six clusters in each model; black rules mark the two predefined communities.
Agent indices align across graph realizations only through those community
memberships, so fine node-level texture in an averaged matrix has no
cross-realization identity interpretation.
The lower panels show equation~\eqref{eq:connected-correlation} by graph
distance, with 95\% complete-cluster bootstrap intervals.  Individual pairs
and updates do not enter the inferential sample size.}{fig:correlations}

Appendix~\ref{app:persistence-binder} and figure~\ref{fig:persistence-binder}
combine the full autocorrelation curve with the fixed-lag integral in
equation~\eqref{eq:autocorrelation}, the Binder statistic in
equation~\eqref{eq:binder}, and empirical magnetization occupancy.  This
combination prevents a Binder value from being interpreted without its
underlying finite-window distribution.  It also separates persistence of the
projected collective coordinate from the block-reversal statistic: the former
is time-symmetric two-point dependence, whereas the latter compares forward
and reversed path words.
The accompanying source table also retains full-, early-half-, and late-half
Binder estimates under cluster-first and pooled-moment rules.  Their
differences diagnose finite-window and aggregation sensitivity rather than
selecting whichever construction yields the larger separation.

At the fixed two-sweep truncation, persistent-minus-Markovized recovery
$\tau_{\mathrm{int}}$ changes are \csname V15QwenTauMemory\endcsname{}
(\csname V15QwenTauMemoryCI\endcsname) updates for Qwen and
\csname V15GraniteTauMemory\endcsname{}
(\csname V15GraniteTauMemoryCI\endcsname) updates for Granite.  The
field-Markovized quench-minus-baseline Binder changes are
\csname V15QwenBinderQuench\endcsname{}
(\csname V15QwenBinderQuenchCI\endcsname) and
\csname V15GraniteBinderQuench\endcsname{}
(\csname V15GraniteBinderQuenchCI\endcsname), respectively.  These estimates
are reported whether positive, negative, or interval-compatible with zero.

\clearpage
\section{Memory and projected temporal asymmetry}
\label{sec:memory-results}

\subsection{Persistent, Markovized, and scrambled histories}
\label{sec:persistent-markovized}

Relative to Markovized state, genuine memory changes adjusted block divergence
by \csname V15HTwo\endcsname{} nats per attempted update (95\% interval
\csname V15HTwoCI\endcsname); the prespecified disposition is
\csname V15HTwoDisposition\endcsname{}.  H2 asks whether history matters
relative to an explicit current-state prompt.

Relative to the scrambled-history
placebo, the change is \csname V15HThree\endcsname{}
(\csname V15HThreeCI\endcsname), disposition
\csname V15HThreeDisposition\endcsname{}.  H3 is more stringent: it asks whether the
temporally related content matters beyond adding a similarly shaped prompt
section.

Several arm-level adjusted values are negative because their observed block
divergence lies below the corresponding finite-sample shuffle-floor mean.  H2
and H3 are paired differences between these untruncated adjusted statistics;
their positive contrasts do not establish positive absolute entropy production
in either arm.

No exact thermodynamic interpretation follows from a positive result.  The
observed block words omit natural-language content and some private state.
Conversely, a null H3 would be informative: prompt presence, bounded memory
depth, model-specific history use, or estimator power could explain the
earlier exploratory and replication-stage contrast without a reproducible
content-specific effect.  Block-length,
pseudocount, conditional-memory-depth, occupancy, and token-balance tables
support this distinction.

Figure~\ref{fig:v15memory} displays the paired cluster contrasts for the two pinned model families.

The prompt-balance diagnostic in figure~\ref{fig:prompt-balance} of
Appendix~\ref{app:sensitivities} records the scope of the format and
approximate-length control.

\resultfigure{figure06_memory_controls.pdf}{0.94\textwidth}{%
Per-cluster genuine-memory contrasts for both pinned model families.  Short
horizontal segments indicate model means.  The scrambled-history arm controls
prompt section and approximate length while destroying temporal relation to
the current trajectory.}{fig:v15memory}

\subsection{Model heterogeneity and estimator dependence}
\label{sec:memory-heterogeneity}

The prospectively pooled contrast retains substantial model heterogeneity.
For persistent minus Markovized state, model-specific means are 0.03041 for
Qwen (five of six cluster effects positive) and 0.07835 for Granite (six of
six positive).  For persistent minus scrambled history they are 0.00689 for
Qwen (four of six positive) and 0.09001 for Granite (six of six positive).
The model-specific Qwen scrambled-control sign-flip value is 0.21875 and is
not separately confirmatory; the prespecified pooled H3 remains supported because
model-stratified cluster pairs were the prespecified units.  This decomposition
supports cross-family direction for the Markovized contrast but not a claim
that content-specific memory effects are uniform across model families.

Estimator sensitivity further limits the scale claim.  The pooled memory
contrasts remain positive over most prespecified block-length and pseudocount
combinations, but their magnitude varies substantially; at block length four
and pseudocount one, persistent-minus-scrambled is only 0.00024 nats per
update and the Granite component is slightly negative.  Conditional-memory
information is likewise not monotonically larger in the persistent arm.
Accordingly, the robust statement concerns the prespecified coarse graining, its
bias floor, and paired controls rather than an estimator-independent entropy-
production rate.

The block-length, pseudocount, and shuffled-floor sensitivities are displayed
in figure~\ref{fig:path-sensitivity} of Appendix~\ref{app:sensitivities}.

Taken together, H1--H4 address distinct claims rather than a single composite
effect: H1 tests field-driven departure, H2 tests persistent history relative
to current-state prompting, H3 tests temporally related content relative to a
format-matched control, and H4 tests fixed-window restoration.  Figure
\ref{fig:effects} places the four prespecified effects and their
cluster-bootstrap intervals in one synthesis while keeping their units
separate.  Under the prespecified tests and multiplicity allocation, all four
criteria are met; the model-specific and estimator boundaries above remain
part of that interpretation.

\resultfigure{figure07_confirmatory_effects.pdf}{0.94\textwidth}{%
Prespecified effects and 95\% cluster-bootstrap intervals.  Distance and
information units occupy separate panels; no claims forest places incompatible
quantities on one numerical axis.  Dispositions use exact cluster sign flips
and the prespecified multiplicity allocation.}{fig:effects}

\section{Reduced dynamical descriptions and closure limits}
\label{sec:closure}

\subsection{Information retained beyond mean order and the kinetic surrogate}
\label{sec:representation}

Magnetization records directional order but can remain near zero when half the
population reorganizes in opposing communities.  Configuration entropy then
records collective state diversity; marginal entropy records local occupancy;
total correlation distinguishes independent uncertainty from coordinated
dependence; effective energy records compatibility with the symmetric
reference; susceptibility and energy variance record fluctuations;
autocorrelation records finite-window persistence; and path reversal records
temporal ordering.  These
quantities are complementary only when their estimators and scales are kept
distinct.

The representation analysis of the earlier quench experiment uses transparent multinomial logistic
regression with leave-one-cluster-out preprocessing.  The delayed permutation
audit tests both absolute balanced accuracy and the increment over order-only
features.  This is not a claim that entropy outperforms all ordinary summaries.
It tests whether a predefined collective representation retains condition
information discarded by a reduced representation.  The field-quench
trajectory supplies a mechanistic example: energy, dependence, entropy rate,
and fluctuation can change while mean magnetization remains balanced.

We fit logistic belief and action responses only to the immutable
microscopic-response table from the prospective replication stage.  The belief model includes private field, delivered
neighbor field times controlled coupling, previous belief, and previous action;
the action model includes updated belief and previous action.  No coefficient
is refitted to either later quench trajectory set.  The surrogate is run on the
same graph generator, initial-state generator, asynchronous update tape, field
schedule, and intervention times.  It is a deliberately low-dimensional
kinetic-Ising comparison rather than a fitted description of every asymmetric
or history-dependent response \cite{aguilera2021}.

Figure~\ref{fig:surrogate} compares the direct trajectories from the earlier
quench experiment with the out-of-sample kinetic surrogate.

\resultfigure{figure08_direct_surrogate_quench.pdf}{0.94\textwidth}{%
Out-of-sample field-quench comparison between direct Qwen trajectories and
the replication-stage kinetic surrogate.  Matching direction, peak timing, normalized
amplitude, recovery, and route area are assessed separately.  Surrogate
failure is evidence about the limits of the closure, not a failed LLM
experiment.}{fig:surrogate}

In the five-coordinate comparison geometry, the direct field-minus-nominal
peak-response difference is only 0.142 units on average and is nonzero in one
of six clusters, whereas the surrogate difference is 1.467 units and is
positive in four clusters.  The corresponding field-minus-nominal
energy--entropy route-area differences are 0.043 and 0.695.  The direct shared
response peaks at the quench boundary in one cluster and is otherwise dominated
by initial relaxation; the surrogate places its peak at the quench or
counter-quench boundary in four clusters.  Both descriptions return close to
their own final shared-response baseline, but the surrogate neither reproduces
the direct amplitude nor its timing.  This also explains why the richer
macrostate distance used in the earlier quench analysis can show a large quench pulse that a short kinetic closure
over a small shared coordinate set misses.

\subsection{Captured response, failure modes, and finite-size context}
\label{sec:closure-limits}

This comparison identifies three possible failure modes.  First, a homogeneous
logistic neighbor coefficient cannot capture state-dependent semantic use.
Second, action and belief persistence can introduce lags not represented by a
single synchronous mean field.  Third, hidden memory and workload can produce
different future laws at the same projected spin configuration.  Adding
unconstrained parameters after observing quench trajectories would obscure
these limits, so the closure remains low-dimensional and out of sample.

CPU trajectories at $N\in\{8,16,32,64\}$ supply a dense effective-model size
context.  They are not direct LLM finite-size scaling and cannot establish a
critical exponent or thermodynamic-limit phase transition.  The direct
$N=16$ paths remain empirical anchors.

The effective-model size context is displayed in figure~\ref{fig:surrogate-size}
of Appendix~\ref{app:sensitivities}.

\section{Discussion}
\label{sec:discussion}

Treating the LLM-agent network as a stochastic interacting system begins with
implemented rather than metaphorical objects: agent instances receive local
prompts, exchange packets, update on a recorded schedule, and change
categorical state.  Magnetization, Shannon entropy, mutual information, total
correlation, and path divergence are then mathematical observables of those
records.  State separation is central to this interpretation because it
distinguishes the complete augmented state that determines a transition from
the local information available to an agent and from the lower-dimensional
projection available for population-level measurement.

The quench experiments use this formulation to measure finite-horizon response
and recovery rather than equilibrium relaxation.  Both model families depart
from their within-model nominal geometries after field reversal, and the
fixed-window restoration contrast supports motion back toward the restored
nominal regime.  The difference between complete Qwen threshold re-entry and
slower or incomplete Granite re-entry within the observation horizon remains
important: the common directional response does not imply a common recovery
time or complete return.  Recovery times and threshold crossings are
trajectory descriptors, and the quench and counter-quench paths need not
retrace because messages, actions, and history change the state reached at
restoration.

The history interventions clarify what memory means under coarse-grained
observation.  Persistent history differs from an explicit current-state
prompt, and comparison with scrambled history asks whether temporally related
content matters beyond prompt format and approximate length.  The distinction
between complete augmented state, observable projection, and rolling
macrostate explains how private memory can generate history dependence in the
projected process.  Positive paired contrasts nevertheless do not establish
positive absolute entropy production: some arm-level bias-adjusted values are
negative, the observed path words omit natural-language content and private
state, and the measured reversal divergence is a coarse-grained projected
statistic rather than an exact physical entropy-production rate.

Model heterogeneity and negative findings limit generalization.  The
persistent-versus-Markovized contrast has the same direction across the two
families, whereas the content-specific scrambled-history contrast is much
weaker and less uniform for Qwen than for Granite.  Estimator choices also
change the magnitude of the memory contrasts.  In the exploratory evidence,
degree- and traffic-matched nonreciprocity did not produce a reliable increase
in irreversibility, and partition and message-corruption trajectories in the
earlier quench experiment remained close to nominal under their tested
construction.  These results prevent the broader assertion that any directed
graph, entropy statistic, or history prompt produces the same collective
effect.

The secondary spatial, persistence, Binder, and closure analyses add distinct
but limited information.  Connected correlations describe how finite-window
organization changes across the fixed modular graph, while autocorrelation
summarizes persistence of the projected collective coordinate.  Binder values
are shown with occupancy and aggregation sensitivities; they do not establish
a crossing, critical slowing down, or a phase transition.  The out-of-sample
kinetic surrogate captures some direction and recovery features but misses
parts of the direct response amplitude and timing.  Its inexpensive size
context is not direct-LLM finite-size scaling, so surrogate failure identifies
limits of the reduced closure rather than a failed agent experiment.

These measurements support an effective statistical-mechanical description,
not a literal identification of the network with a thermodynamic material.
The reference energy is effective because the LLM transition rule is not
derived from its Hamiltonian, and decoding temperature is a generation setting
rather than an independently inferred thermodynamic temperature.  No
free-energy-like coordinate is emphasized because such a temperature is
unavailable.  The value of the framework is instead that it keeps collective
organization, fluctuations, dependence, response, and temporal ordering
operationally distinct while exposing the assumptions behind each reduction.

Several limitations bound the scope.  Two 7B instruction-model families do
not establish universality, and the direct systems contain only $N=16$ agents,
so no thermodynamic-limit transition is claimed.  The binary projection
compresses semantic messages and private state.  The confirmatory study uses
one reciprocal modular topology and one coupling/noise anchor.  The
scrambled-history placebo matches format and approximate token distribution
but cannot make semantic content exchangeable turn by turn.  Short-window
entropy, dependence, autocorrelation, and path estimates remain sensitive to
coarse graining and finite sampling even after their null-floor and
aggregation checks.

Further work therefore requires larger direct systems, additional network
topologies, more model families, and longer stationary trajectories.  Richer
observable-state representations could test which hidden variables account
for projected memory, while prospectively chosen alternative coarse grainings
could establish whether the present estimator dependence persists.  Such work
is needed before thermodynamic-limit, critical, or exact
stochastic-thermodynamic claims would be warranted.

\section{Conclusions}
\label{sec:conclusions}

State-separated LLM-agent instances form a measurable interacting stochastic
system whose collective trajectories admit a disciplined reduced description.
Controlled field quench and restoration create a common test of
order, uncertainty, dependence, effective compatibility, fluctuation, and
temporal response.  The prospective cross-model experiment adds an independent
model family and a format-matched history control to the exploratory and
earlier quench evidence.  The pooled prospective memory effects are positive,
while the weaker Qwen scrambled-control decomposition, Granite's incomplete
threshold re-entry, and estimator sensitivity remain explicit boundary
results.

The statistical-mechanical contribution is therefore not that entropy makes
agents perform better.  It is a formulation and measurement framework that
connects local language-conditioned transitions to collective finite-size
dynamics, exposes what population averages discard, and makes failures of
simple closures explicit.

\section*{Data and code availability}

Code, configurations, figure source data, aggregate results, and reproduction
scripts are available at \url{https://github.com/mantzaris/ThermoAgent}.  The
reported cross-model experiment ran on an NVIDIA GeForce RTX 4090 and used
approximately \csname V15OriginalGPUHours\endcsname{} measured generation
GPU-hours.  Raw prompts, completions, and full trajectories are not distributed
because of size and potential sensitivity, and the original external raw tree
is no longer available; content-addressed manifests and replay and
aggregate-validation artifacts document the reconstruction.

\section*{Acknowledgements}

OpenAI Codex (GPT-5) was used for manuscript organization and language editing,
code assistance, and literature discovery; the author reviewed all resulting
material and takes responsibility for the manuscript.

\appendix

\section{Study chronology, correction, and inferential details}
\label{app:study-inference}

The earlier quench analysis originally derived each field panel's recovery
threshold from that panel's own baseline, although the protocol called for a
threshold fitted only to the training clusters.  The corrected analysis maps
the leave-one-cluster-out training threshold into the held-out recovery
summary without changing any recorded trajectory.  Figure
\ref{fig:v14recovery} retains every cluster-level recovery path, including the
high-response cluster, under the corrected training-only construction.  The
historical maximum-minus-final recovery statistic remains numerically
reproducible but is not used as directional evidence because its construction
is nonnegative for nearly every nonconstant trajectory.

\resultfigure{figure09_cluster_recovery.pdf}{0.94\textwidth}{%
All field-reversal cluster trajectories from the earlier experiment,
including the high-response cluster.  The corrected audit replaces the
held-out-panel threshold with the training-cluster nominal threshold.  The
historical recovery sign statistic is not used as directional evidence.}{fig:v14recovery}

The delayed audit executes the omitted 10,000 cluster-preserving label
permutations.  All four panels stay together within each cluster, and every
replicate refits imputation, scaling, logistic regression, and leave-one-
cluster-out predictions.  Three-, five-, and seven-sweep nominal geometries
are fitted independently.  Each coordinate and observable family is deleted
in turn, and dependence estimates receive circular-shift null floors.  These
were delayed prespecified sensitivity analyses completed after the primary
outcomes because they had been omitted from the initial implementation, not a
new prospective experiment.

The representation permutation retains the four condition panels in every
cluster and permutes labels only within cluster.  For each of 10,000
replicates, every leave-one-cluster-out fold refits imputation, scaling, and the
multinomial logistic model; test-cluster features do not enter preprocessing.
Figure~\ref{fig:v14audit} records the delayed dependence and full-pipeline
permutation audit.

\resultfigure{figure10_delayed_audit.pdf}{0.94\textwidth}{%
Delayed prespecified dependence and permutation audits.  Raw and
bias-adjusted dependence, window sensitivity, and cluster-preserving nulls are
shown explicitly.  These audits were completed after the primary outcomes
because of implementation omissions.}{fig:v14audit}

For the prospective cross-model tests, the one-sided sign-flip probability
enumerates all $2^n$ sign assignments of the observed cluster effects for
$n\leq20$ paired clusters.  The cluster bootstrap samples complete
graph/environment clusters with replacement and never resamples time points
as independent observations.  H1 receives alpha 0.02; H2--H4 receive a
Holm-adjusted family alpha 0.03.  The allocation was fixed before outcomes
were examined.  Intervals are uncertainty summaries and do not replace the exact
paired tests.

\section{Estimator, control, and closure sensitivities}
\label{app:sensitivities}

Total-correlation nulls independently circular-shift each coordinate by a
nonzero offset where possible.  This retains marginal occupancy and each
coordinate's internal temporal ordering while breaking synchronous dependence.
The adjusted estimate is not truncated.  Alternative rolling windows rebuild
the words, covariance, training threshold, and held-out distance rather than
subsampling a five-sweep output.

For block reversal, time-shuffled floors quantify finite-length occupancy
bias; they do not prove stationarity.  Results at block lengths two, three,
and four expose finite-history sensitivity, while conditional mutual
information at history depths one through three provides a separate diagnostic
of incomplete first-order closure.  Figure~\ref{fig:path-sensitivity} displays
the prespecified estimator-sensitivity grid.

\resultfigure{figure11_path_reversal_sensitivity.pdf}{0.94\textwidth}{%
Block-length, pseudocount, and shuffled-bias-floor sensitivity for
bias-adjusted block reversal divergence.  Memory contrasts can be checked
against estimator choices; observable coarse-graining and finite length remain
limiting.}{fig:path-sensitivity}

The scrambled-history arm controls the presence and approximate length of the
history prompt while destroying its temporal relation to the current
trajectory.  Figure~\ref{fig:prompt-balance} records the token-count diagnostic
for that comparison.  It supports approximate prompt-length and format balance
but does not establish turn-by-turn token identity or semantic exchangeability.

\resultfigure{figure12_prompt_balance.pdf}{0.94\textwidth}{%
Scrambled-history prompt-length control.  Per-cluster mean token counts verify
that prompt length and format are approximately matched; semantic content
cannot be exactly token-matched turn by turn.}{fig:prompt-balance}

Finally, CPU trajectories at several sizes provide context for the reduced
kinetic model.  Figure~\ref{fig:surrogate-size} places the direct-system anchor
beside that effective-model size comparison.  These calculations are not
direct-LLM finite-size scaling and do not support a thermodynamic-limit claim.

\resultfigure{figure13_surrogate_size_context.pdf}{0.94\textwidth}{%
Direct $N=16$ anchors in a denser inexpensive effective-model size context.
These are CPU kinetic-surrogate quench responses, not direct-LLM finite-size
scaling.}{fig:surrogate-size}

\section{Finite-window persistence and Binder diagnostics}
\label{app:persistence-binder}

The autocorrelation and Binder diagnostics defined in
equations~\eqref{eq:autocorrelation} and~\eqref{eq:binder} are summarized in
figure~\ref{fig:persistence-binder}.  Each quantity is computed within one
complete model--graph--environment trajectory before uncertainty is summarized
across the six clusters per model.  The autocorrelation sum uses the fixed
two-sweep truncation, while one- and three-sweep truncations remain
machine-readable sensitivity analyses.  Zero-variance phase series are left
undefined rather than assigned artificial persistence, and uncertainty
resamples complete clusters rather than treating time points as independent.

The Binder statistic is likewise calculated trajectory first and describes
the shape of a finite $N=16$ magnetization distribution.  The source-data table
retains full-phase, early-half, and late-half estimates under both cluster-mean
and pooled-moment constructions; near-zero second-moment cases remain
undefined.  Empirical magnetization occupancy is shown alongside the Binder
values so that aggregation-sensitive cumulants remain connected to the
finite-window distributions from which they are calculated.

The autocorrelation sum is not an equilibrium correlation time, and these
short, potentially nonstationary windows cannot establish critical slowing
down.  The Binder values are not a crossing analysis and do not establish a
critical point or phase transition.  Figure~\ref{fig:persistence-binder}
therefore presents persistence, Binder, and occupancy together rather than as
an additional confirmatory hypothesis family.

\resultfigure{figure14_persistence_binder.pdf}{0.98\textwidth}{%
Finite-window persistence and order-parameter shape.  Autocorrelation curves
and their 95\% intervals compare the nominal late window with the three field
arms using complete clusters; integrated correlation sums show all cluster
values at the fixed two-sweep truncation.  Binder cumulants are paired with
empirical belief-magnetization occupancies so their finite-sample meaning
remains visible.  These diagnostics do not establish critical slowing down or
a phase transition.}{fig:persistence-binder}

\clearpage
\section{Computational implementation and validation}
\label{app:provenance}

An agent receives only its own persistent state, the delivered neighbor
packets available at its scheduled update, and the condition-specific prompt
content defined in Section~\ref{sec:system-protocol}.  Its response changes the
scheduled instance's recorded belief, action, commitment, memory code, signal,
and typed tool consequence; the scheduler validates but does not replace those
fields.  Transition records include the scheduled identity, recipient, packet
bytes, and prompt and output hashes.  Dedicated checks confirm that perturbing
one instance's memory or decision provider leaves unrelated agents' private
state unchanged.

Every experimental arm instantiates a fresh network state from its
prespecified panel seeds.  Model weights and the tokenizer are shared
read-only for throughput, but conversational histories, generation caches,
mutable agent objects, and scientific random-number state are not carried
between arms.  Each generation resets its per-decision seed.  The two pinned
model revisions reported in Section~\ref{sec:cross-model-design} use
Transformers 4.55.4, PyTorch 2.8.0+cu128, bitsandbytes 0.47.0, NF4 double
quantization, and BF16 computation under the common decoding settings given
there.

Independent reconstruction validation covered all
\csname V15Trajectories\endcsname{} trajectories and
\csname V15Decisions\endcsname{} retained decisions, with zero replay
mismatches.  Reconstructed aggregate tables agree with the committed reference
within an explicitly audited binary64 cross-platform tolerance, not a
scientific effect-size tolerance.  The repository links the reported results
to the model revisions, configurations, graph and seed manifests, recorded
control schedules, aggregate tables, figure source data, and reproduction
scripts.  The original external raw tree is unavailable, so historical raw-file
digest identity cannot be checked; this limitation is stated separately from
the replay and aggregate comparisons that can be reproduced.

\bibliographystyle{unsrtnat}
{\small
\bibliography{references}
}

\end{document}
